\documentclass[aspectratio=169,8pt]{beamer}
\usetheme{Madrid}
\usepackage[utf8]{inputenc}
\usepackage[T1]{fontenc}
\usepackage{amsmath,amssymb}
\usepackage{booktabs}
\usepackage{enumitem}
\setlist[enumerate]{label=\Alph*.,leftmargin=*,itemsep=1pt,topsep=2pt}
\setlist[itemize]{leftmargin=*,itemsep=1pt,topsep=2pt}
\setbeamerfont{block body}{size=\tiny}
\setbeamerfont{block title}{size=\scriptsize}
\setbeamerfont{frametitle}{size=\footnotesize}
\setbeamerfont{normal text}{size=\tiny}
\setbeamertemplate{navigation symbols}{}
\setbeamertemplate{itemize item}{\tiny$\bullet$}
\setbeamertemplate{itemize subitem}{\tiny$\circ$}

\title[GARP RAI]{GARP Risk \& AI --- Q351--Q450}
\subtitle{Unofficial Exam Prep}
\author{Unofficial}
\date{\today}

\begin{document}
	
	\begin{frame}
		\titlepage
	\end{frame}
	
	% =================================================================
	\section{Data Governance (continued)}
	% =================================================================
	



	



	\begin{frame}{Q357 --- Data Access Control Models}
		\begin{block}{Question}
			Which access control model uses attributes of the user, the resource, and the environment to make authorisation decisions?
			\begin{enumerate}
				\item Role-Based Access Control (RBAC).
				\item Attribute-Based Access Control (ABAC).
				\item Mandatory Access Control (MAC).
				\item Discretionary Access Control (DAC).
			\end{enumerate}
		\end{block}
		\begin{exampleblock}{Answer: B}\end{exampleblock}
		\begin{block}{Explanation}
			ABAC uses attributes of user, resource, and environment. RBAC uses roles; MAC uses security labels; DAC lets owners set permissions.
		\end{block}
	\end{frame}
	
	\begin{frame}{Q358 --- Data Protection Principles}
		\begin{block}{Question}
			Which of the following is NOT a data protection principle commonly found in privacy regulations?
			\begin{enumerate}
				\item Data minimisation.
				\item Purpose limitation.
				\item Storage limitation.
				\item Unlimited retention.
			\end{enumerate}
		\end{block}
		\begin{exampleblock}{Answer: D}\end{exampleblock}
		\begin{block}{Explanation}
			Modern privacy regulations require storage limitation; data must not be kept longer than necessary. Unlimited retention violates this principle.
		\end{block}
	\end{frame}
	

	
	\begin{frame}{Q360 --- Data Encryption}
		\begin{block}{Question}
			Which is the primary purpose of encrypting sensitive data both at rest and in transit?
			\begin{enumerate}
				\item To increase storage capacity.
				\item To protect confidentiality and integrity against unauthorised access.
				\item To reduce processing time.
				\item To comply with marketing standards.
			\end{enumerate}
		\end{block}
		\begin{exampleblock}{Answer: B}\end{exampleblock}
		\begin{block}{Explanation}
			Encryption protects data confidentiality and integrity. Options A, C, and D are unrelated.
		\end{block}
	\end{frame}
	
	% =================================================================
	\section{Model Governance (continued)}
	% =================================================================
	
	\begin{frame}{Q361 --- Model Governance Policy Content}
		\begin{block}{Question}
			Which of the following is typically part of a model governance policy?
			\begin{enumerate}
				\item The definition of a model and non-model.
				\item The model risk tiering methodology.
				\item Validation and documentation requirements.
				\item All of the above.
			\end{enumerate}
		\end{block}
		\begin{exampleblock}{Answer: D}\end{exampleblock}
		\begin{block}{Explanation}
			Model governance policies define models, tiers, validation, documentation, monitoring, and reporting. All three options are standard components.
		\end{block}
	\end{frame}

	\begin{frame}{Q363 --- Model Inventory Essentials}
		\begin{block}{Question}
			Which item is essential in a model inventory?
			\begin{enumerate}
				\item Only the model name.
				\item Model owner, purpose, risk tier, validation status, and interdependencies.
				\item The model's full source code.
				\item The model's marketing brochure.
			\end{enumerate}
		\end{block}
		\begin{exampleblock}{Answer: B}\end{exampleblock}
		\begin{block}{Explanation}
			The inventory tracks ownership, purpose, risk tier, validation, and dependencies. Options A, C, and D are insufficient or inappropriate.
		\end{block}
	\end{frame}
	
	\begin{frame}{Q364 --- Model Landscape Purpose}
		\begin{block}{Question}
			What does a model landscape primarily provide that a model inventory does not?
			\begin{enumerate}
				\item A list of all models.
				\item A map of interactions and interdependencies, highlighting cascade risk and critical nodes.
				\item Model code.
				\item Model marketing material.
			\end{enumerate}
		\end{block}
		\begin{exampleblock}{Answer: B}\end{exampleblock}
		\begin{block}{Explanation}
			Landscape captures interactions; inventory lists models. Options A, C, and D describe inventory or unrelated content.
		\end{block}
	\end{frame}
	
	\begin{frame}{Q365 --- Risk Tiering Factors}
		\begin{block}{Question}
			Which of the following is a standard factor in AI/ML model risk tiering?
			\begin{enumerate}
				\item Materiality of model use.
				\item Complexity of methodology.
				\item Exposure (financial, regulatory, reputational).
				\item All of the above.
			\end{enumerate}
		\end{block}
		\begin{exampleblock}{Answer: D}\end{exampleblock}
		\begin{block}{Explanation}
			Risk tiering considers materiality, complexity, and exposure. All three are essential inputs.
		\end{block}
	\end{frame}
	
	\begin{frame}{Q366 --- Independent Validation}
		\begin{block}{Question}
			What is the primary purpose of independent model validation?
			\begin{enumerate}
				\item To approve the model without further review.
				\item To critically and objectively assess the model's conceptual soundness, data quality, implementation, and performance, by parties independent of the developers.
				\item To rewrite the model code.
				\item To market the model.
			\end{enumerate}
		\end{block}
		\begin{exampleblock}{Answer: B}\end{exampleblock}
		\begin{block}{Explanation}
			Independent validation provides objective challenge. Options A, C, and D are contrary or unrelated.
		\end{block}
	\end{frame}
	
	\begin{frame}{Q367 --- Validation Timing}
		\begin{block}{Question}
			Which of the following best describes when model validation should occur?
			\begin{enumerate}
				\item Only at model retirement.
				\item Before deployment, periodically during use, and after material changes.
				\item Only after a serious incident.
				\item Once, at design time.
			\end{enumerate}
		\end{block}
		\begin{exampleblock}{Answer: B}\end{exampleblock}
		\begin{block}{Explanation}
			Validation is continuous: initial, periodic, and change-based. Options A, C, and D are too narrow.
		\end{block}
	\end{frame}
	
	\begin{frame}{Q368 --- Use Test}
		\begin{block}{Question}
			The ``use test'' in model validation is best described as:
			\begin{enumerate}
				\item A quantitative accuracy metric.
				\item A qualitative assessment of whether the model is appropriately used for its intended purpose in practice.
				\item A test of model speed.
				\item A code-review technique.
			\end{enumerate}
		\end{block}
		\begin{exampleblock}{Answer: B}\end{exampleblock}
		\begin{block}{Explanation}
			The use test is qualitative and focused on real-world use. Options A, C, and D are different.
		\end{block}
	\end{frame}


	
	\begin{frame}{Q371 --- Model Use Limits}
		\begin{block}{Question}
			Why must model use limits be documented and enforced?
			\begin{enumerate}
				\item Because they are optional.
				\item Because using a model outside its validated domain can produce materially incorrect results and unexpected risk.
				\item Because they speed up training.
				\item Because they reduce storage.
			\end{enumerate}
		\end{block}
		\begin{exampleblock}{Answer: B}\end{exampleblock}
		\begin{block}{Explanation}
			Use limits protect against out-of-scope use. Options A, C, and D are incorrect.
		\end{block}
	\end{frame}
	
	\begin{frame}{Q372 --- Change Management}
		\begin{block}{Question}
			A material change to a production model triggers:
			\begin{enumerate}
				\item Immediate deployment without review.
				\item A change management process including impact assessment, documentation, validation, and approval.
				\item Retirement of the model.
				\item No action.
			\end{enumerate}
		\end{block}
		\begin{exampleblock}{Answer: B}\end{exampleblock}
		\begin{block}{Explanation}
			Material changes go through structured change management. Options A, C, and D violate governance.
		\end{block}
	\end{frame}
	


	

	

	
	\begin{frame}{Q378 --- Model Monitoring}
		\begin{block}{Question}
			Ongoing model monitoring should detect:
			\begin{enumerate}
				\item Only training errors.
				\item Performance degradation, data drift, concept drift, and other changes that could affect model reliability.
				\item Only model speed.
				\item Only storage issues.
			\end{enumerate}
		\end{block}
		\begin{exampleblock}{Answer: B}\end{exampleblock}
		\begin{block}{Explanation}
			Monitoring detects degradation and drift. Options A, C, and D are too narrow.
		\end{block}
	\end{frame}
	
	\begin{frame}{Q379 --- Data Drift vs Concept Drift}
		\begin{block}{Question}
			Which statement correctly distinguishes data drift from concept drift?
			\begin{enumerate}
				\item They are the same.
				\item Data drift = change in input distribution; concept drift = change in relationship between inputs and target.
				\item Data drift = change in labels only.
				\item Concept drift = change in storage.
			\end{enumerate}
		\end{block}
		\begin{exampleblock}{Answer: B}\end{exampleblock}
		\begin{block}{Explanation}
			Data drift: input distribution shift. Concept drift: feature-target relationship shift. Options A, C, and D misstate.
		\end{block}
	\end{frame}
	

	% =================================================================
	\section{AI/ML Model Development and Testing}
	% =================================================================
	
	\begin{frame}{Q381 --- Development Process}
		\begin{block}{Question}
			Which of the following is the FIRST step in the model development process?
			\begin{enumerate}
				\item Deploy the model.
				\item Define objectives and scope.
				\item Retire the model.
				\item Audit the model.
			\end{enumerate}
		\end{block}
		\begin{exampleblock}{Answer: B}\end{exampleblock}
		\begin{block}{Explanation}
			Development starts with clear objectives and scope. Options A, C, and D are subsequent or unrelated stages.
		\end{block}
	\end{frame}
	
	\begin{frame}{Q382 --- Data Collection}
		\begin{block}{Question}
			During data collection and preprocessing, analysts should:
			\begin{enumerate}
				\item Skip documentation.
				\item Identify sources, clean data, handle missing values, and document any exclusions with rationale.
				\item Ignore outliers.
				\item Use all data without inspection.
			\end{enumerate}
		\end{block}
		\begin{exampleblock}{Answer: B}\end{exampleblock}
		\begin{block}{Explanation}
			Data preparation requires cleaning, missing-value handling, and documentation. Options A, C, and D skip essential steps.
		\end{block}
	\end{frame}
	
	\begin{frame}{Q383 --- Feature Engineering}
		\begin{block}{Question}
			Feature engineering typically includes:
			\begin{enumerate}
				\item Combining variables, deriving new features, transformations, and dimensionality reduction.
				\item Only removing features.
				\item Only scaling.
				\item Only imputation.
			\end{enumerate}
		\end{block}
		\begin{exampleblock}{Answer: A}\end{exampleblock}
		\begin{block}{Explanation}
			Feature engineering is broad: transformations, new features, combinations, reductions. Options B, C, and D are partial.
		\end{block}
	\end{frame}
	
	\begin{frame}{Q384 --- Model Selection}
		\begin{block}{Question}
			Model selection should be based on:
			\begin{enumerate}
				\item Personal preference.
				\item A comprehensive evaluation of alternatives, considering performance, interpretability, complexity, and business context.
				\item Speed alone.
				\item Cost alone.
			\end{enumerate}
		\end{block}
		\begin{exampleblock}{Answer: B}\end{exampleblock}
		\begin{block}{Explanation}
			Model selection balances multiple criteria. Options A, C, and D are oversimplified.
		\end{block}
	\end{frame}
	
	\begin{frame}{Q385 --- Training and Calibration}
		\begin{block}{Question}
			During model training, cross-validation is used to:
			\begin{enumerate}
				\item Increase training time.
				\item Estimate generalisation performance and tune hyperparameters.
				\item Skip validation.
				\item Remove features.
			\end{enumerate}
		\end{block}
		\begin{exampleblock}{Answer: B}\end{exampleblock}
		\begin{block}{Explanation}
			Cross-validation estimates generalisation and supports hyperparameter tuning. Options A, C, and D are incorrect.
		\end{block}
	\end{frame}
	
	\begin{frame}{Q386 --- Developer Testing}
		\begin{block}{Question}
			Developer testing before deployment is designed to:
			\begin{enumerate}
				\item Skip validation.
				\item Evaluate performance before production release; ensure accuracy, robustness, and stability.
				\item Replace independent validation.
				\item Reduce code size.
			\end{enumerate}
		\end{block}
		\begin{exampleblock}{Answer: B}\end{exampleblock}
		\begin{block}{Explanation}
			Developer testing precedes and complements independent validation. Options A, C, and D misstate its purpose.
		\end{block}
	\end{frame}
	
	\begin{frame}{Q387 --- Testing Stages}
		\begin{block}{Question}
			Which of the following is NOT a standard testing stage?
			\begin{enumerate}
				\item Unit tests.
				\item Integration tests.
				\item Performance tests.
				\item Marketing tests.
			\end{enumerate}
		\end{block}
		\begin{exampleblock}{Answer: D}\end{exampleblock}
		\begin{block}{Explanation}
			Standard testing: unit, component, integration, system, performance, regression, acceptance, adversarial. Marketing is not a testing stage.
		\end{block}
	\end{frame}
	
	\begin{frame}{Q388 --- Adversarial Testing}
		\begin{block}{Question}
			What is adversarial testing in the context of AI models?
			\begin{enumerate}
				\item Standard unit testing.
				\item Actively probing the model with crafted inputs to discover unintended behaviour or vulnerabilities (red teaming).
				\item Testing hardware.
				\item Testing marketing materials.
			\end{enumerate}
		\end{block}
		\begin{exampleblock}{Answer: B}\end{exampleblock}
		\begin{block}{Explanation}
			Adversarial testing = red teaming with crafted inputs. Options A, C, and D are unrelated.
		\end{block}
	\end{frame}
	
	\begin{frame}{Q389 --- Regression Testing}
		\begin{block}{Question}
			Regression testing ensures:
			\begin{enumerate}
				\item New functionality works as expected.
				\item The model continues to function correctly after changes elsewhere in the system.
				\item Faster performance.
				\item Lower cost.
			\end{enumerate}
		\end{block}
		\begin{exampleblock}{Answer: B}\end{exampleblock}
		\begin{block}{Explanation}
			Regression tests verify existing functionality after modifications. Options A, C, and D describe different goals.
		\end{block}
	\end{frame}
	
	\begin{frame}{Q390 --- Acceptance Testing}
		\begin{block}{Question}
			Acceptance testing determines whether:
			\begin{enumerate}
				\item The model meets agreed-upon requirements and is ready for deployment.
				\item The model is the fastest available.
				\item The model is cheapest.
				\item The code is clean.
			\end{enumerate}
		\end{block}
		\begin{exampleblock}{Answer: A}\end{exampleblock}
		\begin{block}{Explanation}
			Acceptance testing verifies readiness for deployment. Options B, C, and D are not the criterion.
		\end{block}
	\end{frame}
	
	% =================================================================
	\section{Responsible AI Governance}
	% =================================================================
	
	\begin{frame}{Q391 --- Ethical AI Framework Benefits}
		\begin{block}{Question}
			Which is a primary benefit of an AI ethics framework?
			\begin{enumerate}
				\item Increased training speed.
				\item Clear guidance for evaluating ethical dilemmas, enhanced trust, and reduced regulatory non-compliance risk.
				\item Reduced model size.
				\item Eliminated bias.
			\end{enumerate}
		\end{block}
		\begin{exampleblock}{Answer: B}\end{exampleblock}
		\begin{block}{Explanation}
			Ethics frameworks provide guidance and reduce regulatory risk. Options A, C, and D are incorrect or overstated.
		\end{block}
	\end{frame}
	
	\begin{frame}{Q392 --- Responsible AI}
		\begin{block}{Question}
			Which of the following best captures the concept of Responsible AI?
			\begin{enumerate}
				\item Maximising model accuracy without regard for consequences.
				\item Ethical, transparent, accountable AI aligned with societal values and mitigating potential harms.
				\item Deploying AI as fast as possible.
				\item Ignoring regulations.
			\end{enumerate}
		\end{block}
		\begin{exampleblock}{Answer: B}\end{exampleblock}
		\begin{block}{Explanation}
			Responsible AI encompasses ethics, transparency, accountability, and harm mitigation. Options A, C, and D contradict the concept.
		\end{block}
	\end{frame}
	
	\begin{frame}{Q393 --- Governance Challenges: Power Asymmetries}
		\begin{block}{Question}
			Which is a governance challenge for AI?
			\begin{enumerate}
				\item Power asymmetries between large tech firms and regulators.
				\item Lack of interpretability.
				\item Unpredictability of AI systems.
				\item All of the above.
			\end{enumerate}
		\end{block}
		\begin{exampleblock}{Answer: D}\end{exampleblock}
		\begin{block}{Explanation}
			Power asymmetries, opacity, and unpredictability are all governance challenges.
		\end{block}
	\end{frame}
	
	\begin{frame}{Q394 --- Lack of AI Ethics Structures}
		\begin{block}{Question}
			Compared to medical ethics, what is missing from AI ethics?
			\begin{enumerate}
				\item Nothing.
				\item Mature professional structures: licensing, mandatory training, ethics committees, and enforceable codes.
				\item Only training.
				\item Only licensing.
			\end{enumerate}
		\end{block}
		\begin{exampleblock}{Answer: B}\end{exampleblock}
		\begin{block}{Explanation}
			AI ethics lacks the mature professional structures of medical ethics. Options A, C, and D are too narrow.
		\end{block}
	\end{frame}
	
	\begin{frame}{Q395 --- Regulatory Landscape EU}
		\begin{block}{Question}
			The EU AI Act regulates AI based on:
			\begin{enumerate}
				\item A single uniform standard.
				\item A risk-tiered approach: unacceptable, high, limited, minimal.
				\item Self-certification only.
				\item No requirements.
			\end{enumerate}
		\end{block}
		\begin{exampleblock}{Answer: B}\end{exampleblock}
		\begin{block}{Explanation}
			The AI Act uses a risk-tiered approach. Options A, C, and D misstate the structure.
		\end{block}
	\end{frame}
	

	\begin{frame}{Q397 --- Regulatory Landscape China}
		\begin{block}{Question}
			Which best describes China's AI regulation?
			\begin{enumerate}
				\item No AI regulation.
				\item Targeted regulations including the Provisional Administrative Measures for Generative AI Services and PIPL.
				\item Only medical AI regulation.
				\item Complete industry self-regulation.
			\end{enumerate}
		\end{block}
		\begin{exampleblock}{Answer: B}\end{exampleblock}
		\begin{block}{Explanation}
			China has adopted targeted AI regulations. Options A, C, and D are incorrect.
		\end{block}
	\end{frame}
	

	
	\begin{frame}{Q471 --- Fairness} \begin{block}{Question} Fairness within Responsible AI is BEST described as: \begin{enumerate} \item Treating all individuals identically \item Preventing unjustified bias and discrimination \item Maximizing model accuracy \item Maximizing automation \end{enumerate} \end{block} \begin{exampleblock}{Answer: B}\end{exampleblock} \end{frame}
	
	\begin{frame}{Q399 --- Ethical Theory Comparison}
		\begin{block}{Question}
			Which statement correctly distinguishes consequentialism, deontology, and virtue ethics?
			\begin{enumerate}
				\item All three are identical.
				\item Consequentialism focuses on outcomes; deontology on duties; virtue ethics on character.
				\item All three focus only on outcomes.
				\item All three ignore outcomes.
			\end{enumerate}
		\end{block}
		\begin{exampleblock}{Answer: B}\end{exampleblock}
		\begin{block}{Explanation}
			The three frameworks differ in focus. Options A, C, and D misstate them.
		\end{block}
	\end{frame}
	
	\begin{frame}{Q400 --- Justice Conflicts}
		\begin{block}{Question}
			Which statement about justice is correct?
			\begin{enumerate}
				\item Different distributive justice concepts can conflict.
				\item Impartiality is not required for procedural justice.
				\item Autonomy is secondary to social justice.
				\item Meritocracy is not a distributive justice concept.
			\end{enumerate}
		\end{block}
		\begin{exampleblock}{Answer: A}\end{exampleblock}
		\begin{block}{Explanation}
			Distributive justice concepts can conflict. Options B, C, and D are incorrect.
		\end{block}
	\end{frame}
	
	
	
	
	\begin{frame}{Q400 --- Justice Conflicts} \begin{block}{Question} Which statement about justice is correct? \begin{enumerate} \item Different distributive justice concepts can conflict. \item Impartiality is not required for procedural justice. \item Autonomy is secondary to social justice. \item Meritocracy is not a distributive justice concept. \end{enumerate} \end{block} \begin{exampleblock}{Answer: A}\end{exampleblock} \begin{block}{Explanation} Different theories of distributive justice may produce conflicting outcomes. Equality, need, merit, and utility can result in different distributions. \end{block} \end{frame} \begin{frame}{Q401 --- Equality vs Equity} \begin{block}{Question} What best distinguishes equity from equality? \begin{enumerate} \item Equality gives identical outcomes while equity considers relevant differences. \item Equity requires identical treatment. \item Equality always produces fair outcomes. \item Equity ignores historical disadvantage. \end{enumerate} \end{block} \begin{exampleblock}{Answer: A}\end{exampleblock} \begin{block}{Explanation} Equality focuses on equal treatment, while equity may require different treatment to achieve fair opportunities or outcomes. \end{block} \end{frame} \begin{frame}{Q402 --- Procedural Justice} \begin{block}{Question} Which principle is most central to procedural justice? \begin{enumerate} \item Transparency and impartial decision making \item Maximum economic efficiency \item Equal outcomes \item Maximization of shareholder value \end{enumerate} \end{block} \begin{exampleblock}{Answer: A}\end{exampleblock} \begin{block}{Explanation} Procedural justice focuses on whether decisions are made through fair, consistent, transparent, and unbiased processes. \end{block} \end{frame} \begin{frame}{Q403 --- Meritocracy} \begin{block}{Question} Under a merit-based conception of distributive justice, benefits should primarily be allocated according to: \begin{enumerate} \item Individual contribution and achievement \item Random assignment \item Equal outcomes \item Social status \end{enumerate} \end{block} \begin{exampleblock}{Answer: A}\end{exampleblock} \begin{block}{Explanation} Meritocracy allocates rewards based on contribution, skill, effort, or achievement. \end{block} \end{frame} \begin{frame}{Q404 --- Need-Based Justice} \begin{block}{Question} A need-based theory of justice would most likely support: \begin{enumerate} \item Allocating resources according to individual need \item Equal distribution regardless of circumstances \item Allocation based solely on performance \item Allocation by lottery \end{enumerate} \end{block} \begin{exampleblock}{Answer: A}\end{exampleblock} \begin{block}{Explanation} Need-based approaches prioritize those with greater needs rather than equal or merit-based distributions. \end{block} \end{frame} \begin{frame}{Q405 --- Algorithmic Fairness Tradeoffs} \begin{block}{Question} What is a common challenge in algorithmic fairness? \begin{enumerate} \item Multiple fairness metrics may be impossible to satisfy simultaneously. \item Fairness eliminates all model risk. \item Fairness guarantees equal outcomes. \item Transparency removes all bias. \end{enumerate} \end{block} \begin{exampleblock}{Answer: A}\end{exampleblock} \begin{block}{Explanation} Different mathematical definitions of fairness can conflict, requiring governance decisions and tradeoff analysis. \end{block} \end{frame} \begin{frame}{Q406 --- Corrective Justice} \begin{block}{Question} Corrective justice is primarily concerned with: \begin{enumerate} \item Remedying harms or wrongs \item Equal distribution of resources \item Economic growth \item Maximizing efficiency \end{enumerate} \end{block} \begin{exampleblock}{Answer: A}\end{exampleblock} \begin{block}{Explanation} Corrective justice seeks to restore parties after harm, error, discrimination, or other wrongdoing. \end{block} \end{frame} \begin{frame}{Q407 --- Responsible AI and Justice} \begin{block}{Question} Which action best promotes justice in AI systems? \begin{enumerate} \item Assessing disparate impacts across affected groups \item Maximizing accuracy regardless of consequences \item Eliminating human oversight \item Ignoring stakeholder concerns \end{enumerate} \end{block} \begin{exampleblock}{Answer: A}\end{exampleblock} \begin{block}{Explanation} Justice-oriented AI governance evaluates how benefits, burdens, and errors are distributed among stakeholders. \end{block} \end{frame}
	% =================================================================
	\section{GenAI Governance}
	% =================================================================

	
	\begin{frame}{Q402 --- GenAI Hallucination}
		\begin{block}{Question}
			Which is the best mitigation for hallucination risk in a legal-advice GenAI application?
			\begin{enumerate}
				\item Increase temperature.
				\item Ground the model with retrieval-augmented generation (RAG) using authoritative sources, and require human review for critical outputs.
				\item Use a smaller model.
				\item Disable the model.
			\end{enumerate}
		\end{block}
		\begin{exampleblock}{Answer: B}\end{exampleblock}
		\begin{block}{Explanation}
			RAG grounds output; human review catches remaining errors. Options A, C, and D are inadequate or extreme.
		\end{block}
	\end{frame}
	
	
	\begin{frame}{Q404 --- GenAI Governance Levers}
		\begin{block}{Question}
			Which is an internal governance lever for GenAI?
			\begin{enumerate}
				\item Assigning clear ownership and accountability across the GenAI lifecycle.
				\item Tracking external regulatory changes.
				\item Relying only on external standards.
				\item Outsourcing all governance.
			\end{enumerate}
		\end{block}
		\begin{exampleblock}{Answer: A}\end{exampleblock}
		\begin{block}{Explanation}
			Internal levers include ownership, review processes, escalation, training. Options B and C are external; D is contrary to governance.
		\end{block}
	\end{frame}
	
	\begin{frame}{Q405 --- GenAI External Constraints}
		\begin{block}{Question}
			Which is an external governance constraint for GenAI?
			\begin{enumerate}
				\item Evolving regulatory requirements.
				\item Contractual terms with model providers.
				\item Emerging standards (e.g., NIST AI RMF, ISO/IEC 42001).
				\item All of the above.
			\end{enumerate}
		\end{block}
		\begin{exampleblock}{Answer: D}\end{exampleblock}
		\begin{block}{Explanation}
			All three are external governance constraints on GenAI deployments.
		\end{block}
	\end{frame}
	



	\begin{frame}{Q409 --- GenAI Model Inventory}
		\begin{block}{Question}
			When registering a GenAI system in a model inventory, which additional aspects should be captured beyond traditional models?
			\begin{enumerate}
				\item The underlying foundation model, prompt templates, and any retrieval/knowledge sources.
				\item Only the API provider.
				\item Only the temperature parameter.
				\item Only the compute costs.
			\end{enumerate}
		\end{block}
		\begin{exampleblock}{Answer: A}\end{exampleblock}
		\begin{block}{Explanation}
			GenAI inventory requires capturing foundation model, prompts, and retrieval data. Options B, C, and D are partial.
		\end{block}
	\end{frame}
	
	\begin{frame}{Q410 --- GenAI Validation}
		\begin{block}{Question}
			Which is a unique challenge when validating GenAI models?
			\begin{enumerate}
				\item Same challenges as traditional models.
				\item Open-ended output, hallucinations, multi-modality, and difficulty establishing ground truth.
				\item Only speed.
				\item Only cost.
			\end{enumerate}
		\end{block}
		\begin{exampleblock}{Answer: B}\end{exampleblock}
		\begin{block}{Explanation}
			GenAI validation faces open-ended output, hallucination, multi-modality, and ground-truth issues. Options A, C, and D underestimate the challenge.
		\end{block}
	\end{frame}
	
	% =================================================================
	\section{Advanced Governance and Risk Integration}
	% =================================================================
	
	\begin{frame}{Q411 --- Data-First Governance}
		\begin{block}{Question}
			Why does the AI/ML paradigm shift governance emphasis toward data?
			\begin{enumerate}
				\item Data is irrelevant to AI.
				\item AI/ML models often learn directly from data, making data quality, provenance, and bias central to model behaviour and risk.
				\item Models do not depend on data.
				\item Governance is unchanged.
			\end{enumerate}
		\end{block}
		\begin{exampleblock}{Answer: B}\end{exampleblock}
		\begin{block}{Explanation}
			AI/ML models learn from data; data governance is therefore central. Options A, C, and D are incorrect.
		\end{block}
	\end{frame}
	
	\begin{frame}{Q412 --- Ethical AI in Data}
		\begin{block}{Question}
			Which attribute is most relevant when classifying AI/ML models for ethical review?
			\begin{enumerate}
				\item Model accuracy only.
				\item Data bias, explainability, and recalibration needs, among others.
				\item Model speed.
				\item Compute cost.
			\end{enumerate}
		\end{block}
		\begin{exampleblock}{Answer: B}\end{exampleblock}
		\begin{block}{Explanation}
			Ethical review focuses on bias, explainability, and recalibration. Options A, C, and D are secondary.
		\end{block}
	\end{frame}
	
	\begin{frame}{Q413 --- Model Landscape Governance}
		\begin{block}{Question}
			A model landscape captures interdependencies because:
			\begin{enumerate}
				\item Models operate in isolation.
				\item Outputs of one model are often inputs to others, so failures can cascade across an organisation.
				\item Dependencies are irrelevant.
				\item Only audit cares about them.
			\end{enumerate}
		\end{block}
		\begin{exampleblock}{Answer: B}\end{exampleblock}
		\begin{block}{Explanation}
			Interdependencies create cascading risk. Options A, C, and D are incorrect.
		\end{block}
	\end{frame}
	
	\begin{frame}{Q414 --- Aggregate Model Risk}
		\begin{block}{Question}
			Aggregate model risk can be influenced by:
			\begin{enumerate}
				\item Only individual model quality.
				\item Model interdependencies, shared assumptions, and correlated methodology.
				\item Only data size.
				\item Only IT infrastructure.
			\end{enumerate}
		\end{block}
		\begin{exampleblock}{Answer: B}\end{exampleblock}
		\begin{block}{Explanation}
			Aggregate risk arises from interdependencies and shared assumptions. Options A, C, and D miss this.
		\end{block}
	\end{frame}
	
	\begin{frame}{Q415 --- Data Aggregation Risk}
		\begin{block}{Question}
			Why can aggregating data across sources introduce model risk?
			\begin{enumerate}
				\item Aggregation always improves quality.
				\item Inconsistent definitions, timeframes, or measurement scales can distort features and lead to biased models.
				\item Aggregation eliminates risk.
				\item Aggregation is always safe.
			\end{enumerate}
		\end{block}
		\begin{exampleblock}{Answer: B}\end{exampleblock}
		\begin{block}{Explanation}
			Aggregation can introduce inconsistencies. Options A, C, and D are incorrect.
		\end{block}
	\end{frame}
	
	\begin{frame}{Q416 --- Feedback Loops in Production}
		\begin{block}{Question}
			Which is an example of a production feedback loop creating model risk?
			\begin{enumerate}
				\item A model trained on past decisions then used to make new decisions, which become the next training data.
				\item A model without historical data.
				\item A model that is never deployed.
				\item A model that is retired.
			\end{enumerate}
		\end{block}
		\begin{exampleblock}{Answer: A}\end{exampleblock}
		\begin{block}{Explanation}
			Feedback loops reinforce biases. Options B, C, and D are not feedback loops.
		\end{block}
	\end{frame}
	
	\begin{frame}{Q417 --- Data and Model Lineage}
		\begin{block}{Question}
			Why is data and model lineage important for AI governance?
			\begin{enumerate}
				\item To track data and model transformations end-to-end, enabling reproducibility, auditability, and root-cause analysis.
				\item To speed up training.
				\item To reduce costs.
				\item To remove features.
			\end{enumerate}
		\end{block}
		\begin{exampleblock}{Answer: A}\end{exampleblock}
		\begin{block}{Explanation}
			Lineage supports reproducibility and audit. Options B, C, and D are secondary.
		\end{block}
	\end{frame}
	
	\begin{frame}{Q418 --- Data Stewardship in AI}
		\begin{block}{Question}
			Data stewards are responsible for:
			\begin{enumerate}
				\item Only fixing typos.
				\item Managing data quality, access, definitions, and use within a governance framework.
				\item Building models.
				\item Marketing data.
			\end{enumerate}
		\end{block}
		\begin{exampleblock}{Answer: B}\end{exampleblock}
		\begin{block}{Explanation}
			Data stewards manage quality, access, and definitions. Options A, C, and D are too narrow or unrelated.
		\end{block}
	\end{frame}
	
	\begin{frame}{Q419 --- Data Custodian Role}
		\begin{block}{Question}
			Data custodians are primarily responsible for:
			\begin{enumerate}
				\item Business ownership of data.
				\item Technical storage, security, and access controls for data.
				\item Building models.
				\item Marketing data.
			\end{enumerate}
		\end{block}
		\begin{exampleblock}{Answer: B}\end{exampleblock}
		\begin{block}{Explanation}
			Custodians handle technical storage and access. Stewards handle business data quality; owners provide governance. Options A, C, and D are not custody roles.
		\end{block}
	\end{frame}
	
	\begin{frame}{Q420 --- Data Governance Reporting}
		\begin{block}{Question}
			Data governance reporting to the board typically includes:
			\begin{enumerate}
				\item Only storage costs.
				\item Data quality metrics, incidents, compliance status, and key risks.
				\item Only IT staff lists.
				\item Only marketing spend.
			\end{enumerate}
		\end{block}
		\begin{exampleblock}{Answer: B}\end{exampleblock}
		\begin{block}{Explanation}
			Board-level reporting covers quality, incidents, compliance, and risk. Options A, C, and D are irrelevant.
		\end{block}
	\end{frame}
	
	% =================================================================
	\section{Advanced Model Risk Topics}
	% =================================================================
	
	\begin{frame}{Q421 --- Model Revalidation Trigger}
		\begin{block}{Question}
			Which of the following is a trigger for ad-hoc model revalidation?
			\begin{enumerate}
				\item Only scheduled review.
				\item Significant changes in inputs, environment, or performance deviations.
				\item Only cost increases.
				\item Only headcount changes.
			\end{enumerate}
		\end{block}
		\begin{exampleblock}{Answer: B}\end{exampleblock}
		\begin{block}{Explanation}
			Ad-hoc revalidation is triggered by material changes or performance issues. Options A, C, and D are not standard triggers.
		\end{block}
	\end{frame}
	
	
	
	\begin{frame}{Q423 --- Recalibration vs Redevelopment}
		\begin{block}{Question}
			Recalibration of a model is appropriate when:
			\begin{enumerate}
				\item The model's structure is wrong.
				\item The core structure remains valid but parameters need updating in light of new data.
				\item The model is being retired.
				\item The model has never been used.
			\end{enumerate}
		\end{block}
		\begin{exampleblock}{Answer: B}\end{exampleblock}
		\begin{block}{Explanation}
			Recalibration updates parameters; structural issues require redevelopment. Options A, C, and D describe different actions.
		\end{block}
	\end{frame}
	
	\begin{frame}{Q424 --- Model Risk Appetite Statement Content}
		\begin{block}{Question}
			A model risk appetite statement typically includes:
			\begin{enumerate}
				\item Detailed code.
				\item The organisation's tolerance for model risk, escalation triggers, and consequences.
				\item Only model inventory.
				\item Only validation reports.
			\end{enumerate}
		\end{block}
		\begin{exampleblock}{Answer: B}\end{exampleblock}
		\begin{block}{Explanation}
			Appetite statements define tolerance and triggers. Options A, C, and D are unrelated.
		\end{block}
	\end{frame}
	
	\begin{frame}{Q425 --- Model Risk Taxonomy}
		\begin{block}{Question}
			Which of the following is NOT typically included in a model risk taxonomy?
			\begin{enumerate}
				\item Conceptual risk.
				\item Data risk.
				\item Implementation risk.
				\item Employee satisfaction risk.
			\end{enumerate}
		\end{block}
		\begin{exampleblock}{Answer: D}\end{exampleblock}
		\begin{block}{Explanation}
			Model risk taxonomies focus on conceptual, data, and implementation risk (plus usage risk). Employee satisfaction is not standard.
		\end{block}
	\end{frame}
	
	\begin{frame}{Q426 --- Model Risk Register}
		\begin{block}{Question}
			A model risk register typically tracks:
			\begin{enumerate}
				\item Only model code.
				\item Identified risks, owners, mitigation actions, deadlines, and status.
				\item Only training data.
				\item Only IT systems.
			\end{enumerate}
		\end{block}
		\begin{exampleblock}{Answer: B}\end{exampleblock}
		\begin{block}{Explanation}
			A risk register tracks identified risks, owners, actions, deadlines, and status. Options A, C, and D are unrelated.
		\end{block}
	\end{frame}
	
	\begin{frame}{Q427 --- Model Risk Committee}
		\begin{block}{Question}
			A model risk committee is typically responsible for:
			\begin{enumerate}
				\item Approving models and overseeing model risk governance at the enterprise level.
				\item Entering data.
				\item Training models.
				\item Marketing models.
			\end{enumerate}
		\end{block}
		\begin{exampleblock}{Answer: A}\end{exampleblock}
		\begin{block}{Explanation}
			The model risk committee provides enterprise oversight and approvals. Options B, C, and D are operational roles.
		\end{block}
	\end{frame}
	
	\begin{frame}{Q428 --- Model Risk Training}
		\begin{block}{Question}
			Effective model risk training programs should include:
			\begin{enumerate}
				\item Only technical content.
				\item Technical content plus governance, ethics, and regulatory expectations.
				\item Only ethics.
				\item Only regulation.
			\end{enumerate}
		\end{block}
		\begin{exampleblock}{Answer: B}\end{exampleblock}
		\begin{block}{Explanation}
			Model risk training should be comprehensive: technical, governance, ethics, regulation.
		\end{block}
	\end{frame}
	
	\begin{frame}{Q429 --- Model Risk Culture}
		\begin{block}{Question}
			A strong model risk culture is characterised by:
			\begin{enumerate}
				\item Silence about risks.
				\item Awareness, accountability, open communication, and continuous learning around model risk.
				\item Only senior leaders caring about risk.
				\item Ignoring model limitations.
			\end{enumerate}
		\end{block}
		\begin{exampleblock}{Answer: B}\end{exampleblock}
		\begin{block}{Explanation}
			A strong culture values awareness, accountability, openness, and learning. Options A, C, and D contradict this.
		\end{block}
	\end{frame}
	
	\begin{frame}{Q430 --- Model Risk Reporting Frequency}
		\begin{block}{Question}
			Model risk reporting to governance bodies typically occurs:
			\begin{enumerate}
				\item Only when problems arise.
				\item On a scheduled basis (e.g., quarterly) and ad hoc for significant issues.
				\item Never.
				\item Only at year-end.
			\end{enumerate}
		\end{block}
		\begin{exampleblock}{Answer: B}\end{exampleblock}
		\begin{block}{Explanation}
			Reporting is scheduled and ad hoc as needed. Options A, C, and D are insufficient.
		\end{block}
	\end{frame}
	
	% =================================================================
	\section{Ethics and Privacy (Extended)}
	% =================================================================
	
	\begin{frame}{Q431 --- GDPR Legal Basis}
		\begin{block}{Question}
			Which of the following is a lawful basis for processing personal data under GDPR?
			\begin{enumerate}
				\item Consent.
				\item Contract performance.
				\item Legal obligation.
				\item All of the above.
			\end{enumerate}
		\end{block}
		\begin{exampleblock}{Answer: D}\end{exampleblock}
		\begin{block}{Explanation}
			GDPR Art. 6 specifies lawful bases: consent, contract, legal obligation, vital interests, public task, legitimate interests. Options A, B, and C are each valid bases.
		\end{block}
	\end{frame}
	
	\begin{frame}{Q432 --- Data Subject Rights}
		\begin{block}{Question}
			Which is a data subject right under GDPR?
			\begin{enumerate}
				\item Right to access.
				\item Right to rectification.
				\item Right to erasure.
				\item All of the above.
			\end{enumerate}
		\end{block}
		\begin{exampleblock}{Answer: D}\end{exampleblock}
		\begin{block}{Explanation}
			GDPR grants multiple rights: access, rectification, erasure, restriction, portability, objection. Options A, B, and C are each among them.
		\end{block}
	\end{frame}
	
	\begin{frame}{Q433 --- Privacy by Design}
		\begin{block}{Question}
			Privacy by design requires:
			\begin{enumerate}
				\item Adding privacy features at the end.
				\item Embedding privacy protections into the design and architecture of systems from the outset.
				\item Ignoring privacy until a complaint.
				\item Only privacy policies.
			\end{enumerate}
		\end{block}
		\begin{exampleblock}{Answer: B}\end{exampleblock}
		\begin{block}{Explanation}
			Privacy by design embeds protections from the start. Options A, C, and D are inconsistent with the principle.
		\end{block}
	\end{frame}
	
	\begin{frame}{Q434 --- Anonymisation vs Pseudonymisation}
		\begin{block}{Question}
			What is the difference between anonymisation and pseudonymisation?
			\begin{enumerate}
				\item They are the same.
				\item Anonymisation irreversibly removes identifiability; pseudonymisation replaces identifiers with a key that allows re-identification.
				\item Pseudonymisation is irreversible.
				\item Anonymisation always fails.
			\end{enumerate}
		\end{block}
		\begin{exampleblock}{Answer: B}\end{exampleblock}
		\begin{block}{Explanation}
			Anonymisation is irreversible; pseudonymisation can be reversed with the key. Options A, C, and D misstate the distinction.
		\end{block}
	\end{frame}
	
	\begin{frame}{Q435 --- Differential Privacy}
		\begin{block}{Question}
			What is the purpose of differential privacy?
			\begin{enumerate}
				\item To remove all data.
				\item To allow statistical analysis of a dataset while protecting individual records, by adding calibrated noise.
				\item To encrypt the data.
				\item To compress data.
			\end{enumerate}
		\end{block}
		\begin{exampleblock}{Answer: B}\end{exampleblock}
		\begin{block}{Explanation}
			Differential privacy adds noise to protect individual records while preserving aggregate statistics. Options A, C, and D are different techniques.
		\end{block}
	\end{frame}
	
	\begin{frame}{Q436 --- Ethical AI Review}
		\begin{block}{Question}
			Which of the following is typically part of an ethical AI review?
			\begin{enumerate}
				\item Only model accuracy.
				\item Assessing fairness, transparency, accountability, and potential harms.
				\item Only speed.
				\item Only cost.
			\end{enumerate}
		\end{block}
		\begin{exampleblock}{Answer: B}\end{exampleblock}
		\begin{block}{Explanation}
			Ethical AI reviews assess fairness, transparency, accountability, and harms. Options A, C, and D are insufficient.
		\end{block}
	\end{frame}
	
n
	
	\begin{frame}{Q440 --- Algorithmic Impact Assessment}
		\begin{block}{Question}
			An algorithmic impact assessment (AIA) typically evaluates:
			\begin{enumerate}
				\item Only model accuracy.
				\item Risks and impacts on individuals, groups, and society, especially for high-risk use cases.
				\item Only cost.
				\item Only speed.
			\end{enumerate}
		\end{block}
		\begin{exampleblock}{Answer: B}\end{exampleblock}
		\begin{block}{Explanation}
			AIAs evaluate broader risks and impacts. Options A, C, and D are too narrow.
		\end{block}
	\end{frame}
	
	% =================================================================
	\section{Miscellaneous Advanced Topics}
	% =================================================================

	
	\begin{frame}{Q442 --- Feature Importance Limitations}
		\begin{block}{Question}
			Which is a limitation of impurity-based feature importance for tree ensembles?
			\begin{enumerate}
				\item It is always accurate.
				\item It can be biased toward high-cardinality or correlated features.
				\item It cannot be computed.
				\item It replaces SHAP.
			\end{enumerate}
		\end{block}
		\begin{exampleblock}{Answer: B}\end{exampleblock}
		\begin{block}{Explanation}
			Impurity-based importance is biased; permutation importance or SHAP are alternatives. Options A, C, and D are incorrect.
		\end{block}
	\end{frame}
	
	\begin{frame}{Q443 --- Partial Dependence Plots}
		\begin{block}{Question}
			Partial dependence plots show:
			\begin{enumerate}
				\item Raw data.
				\item The marginal effect of a feature on the model's prediction, averaged over other features.
				\item Model code.
				\item Only training errors.
			\end{enumerate}
		\end{block}
		\begin{exampleblock}{Answer: B}\end{exampleblock}
		\begin{block}{Explanation}
			PDPs display marginal feature effects. Options A, C, and D are unrelated.
		\end{block}
	\end{frame}
	
	\begin{frame}{Q444 --- SHAP Consistency}
		\begin{block}{Question}
			Which property does SHAP have that many other feature attribution methods lack?
			\begin{enumerate}
				\item Non-determinism.
				\item Consistency --- if a feature's contribution increases, its SHAP value does not decrease.
				\item Randomness.
				\item Irreproducibility.
			\end{enumerate}
		\end{block}
		\begin{exampleblock}{Answer: B}\end{exampleblock}
		\begin{block}{Explanation}
			SHAP satisfies consistency, along with local accuracy and missingness. Options A, C, and D are not properties of SHAP.
		\end{block}
	\end{frame}
	
	\begin{frame}{Q445 --- LIME Limitations}
		\begin{block}{Question}
			Which is a limitation of LIME?
			\begin{enumerate}
				\item It is always globally faithful.
				\item It provides local approximations that may not be stable across similar instances and are not guaranteed to be globally faithful.
				\item It has no applications.
				\item It replaces model training.
			\end{enumerate}
		\end{block}
		\begin{exampleblock}{Answer: B}\end{exampleblock}
		\begin{block}{Explanation}
			LIME is local and can be unstable. Options A, C, and D misstate.
		\end{block}
	\end{frame}
	
	\begin{frame}{Q446 --- Counterfactual Explanations}
		\begin{block}{Question}
			Counterfactual explanations are particularly useful because:
			\begin{enumerate}
				\item They describe model code.
				\item They provide actionable, human-understandable explanations of what would need to change for a different decision.
				\item They replace SHAP.
				\item They measure accuracy.
			\end{enumerate}
		\end{block}
		\begin{exampleblock}{Answer: B}\end{exampleblock}
		\begin{block}{Explanation}
			Counterfactuals provide actionable insights. Options A, C, and D are incorrect.
		\end{block}
	\end{frame}


	
	\begin{frame}{Q449 --- Model Drift Response}
		\begin{block}{Question}
			When data drift is detected in a production model, which of the following is generally the appropriate response?
			\begin{enumerate}
				\item Ignore it.
				\item Investigate cause, consider recalibration or retraining, and update governance documentation accordingly.
				\item Delete the model.
				\item Change the target.
			\end{enumerate}
		\end{block}
		\begin{exampleblock}{Answer: B}\end{exampleblock}
		\begin{block}{Explanation}
			Drift triggers investigation and possible recalibration/retraining. Options A, C, and D are not standard.
		\end{block}
	\end{frame}
	
	\begin{frame}{Q450 --- Responsible AI Closing Principle}
		\begin{block}{Question}
			Which of the following best summarises the guiding principle of Responsible AI governance in financial services?
			\begin{enumerate}
				\item Maximise model complexity.
				\item Balance performance with governance, validation, monitoring, fairness, explainability, and ethical alignment with regulatory and business requirements.
				\item Deploy all models without validation.
				\item Ignore risks.
			\end{enumerate}
		\end{block}
		\begin{exampleblock}{Answer: B}\end{exampleblock}
		\begin{block}{Explanation}
			Responsible AI balances performance with governance, validation, monitoring, fairness, explainability, and ethics. Options A, C, and D contradict the principle.
		\end{block}
	\end{frame}
	
	
\end{document}